\ifdefined\gkver\else\def\gkver{student}\fi
\documentclass[\gkver]{gaokaozhenti}
\begin{document}

\chapter{2005年广东理综}
%% paperType: 理综物理部分

\begin{enumerate}
\item
%% number: 1
%% paperName: 2005年广东理综第 1 题 4 分
%% typeId: 2
%% type: 多选题
%% chapter: 运动和力
%% point: 牛顿第一定律
%% method: 无
%% score: 4
%% degree: 650
%% duplicateId: 0
%% body:
一汽车在路面情况相同的公路上直线行驶，下面关于车速、惯性、质量和滑行路程的讨论，正确的是\xzanswer{BC}

\fourchoices[answer=BC]
{车速越大，它的惯性越大}
{质量越大，它的惯性越大}
{车速越大，刹车后滑行的路程越长}
{车速越大，刹车后滑行的路程越长，所以惯性越大}

%% answer: BC
%% memo:
\memoanswer{
本题原卷及材料中未提供详解，待补充。
}

\item
%% number: 2
%% paperName: 2005年广东理综第 2 题 4 分
%% typeId: 1
%% type: 单选题
%% chapter: 原子和原子核
%% point: 核裂变
%% method: 无
%% score: 4
%% degree: 650
%% duplicateId: 0
%% body:
下列说法不正确的是\xzanswer{D}

\fourchoices[answer=D]
{$\ce{^{2}_{1}H + ^{3}_{1}H -> ^{4}_{2}He + ^{1}_{0}n}$ 是聚变}
{$\ce{^{235}_{92}U + ^{1}_{0}n -> ^{140}_{54}Xe + ^{94}_{38}Sr + 2^{1}_{0}n}$ 是裂变}
{$\ce{^{226}_{88}Ra -> ^{222}_{86}Rn + ^{4}_{2}He}$ 是 $\alpha$ 衰变}
{$\ce{^{24}_{11}Na -> ^{24}_{12}Mg + ^{0}_{-1}e}$ 是裂变}

%% answer: D
%% memo:
\memoanswer{
本题原卷及材料中未提供详解，待补充。
}

\item
%% number: 3
%% paperName: 2005年广东理综第 3 题 4 分
%% typeId: 1
%% type: 单选题
%% chapter: 机械波
%% point: 波的多解性
%% method: 无
%% score: 4
%% degree: 650
%% duplicateId: 0
%% body:
一列沿 $x$ 轴正方向传播的简谐横波，$t=0$ 时刻的波形如图中实线所示，$t=0.2\Us$ 时刻的波形如图中的虚线所示，则\xzanswer{C}

\onepicture[w=8.5cm]{figs/03.png}

\fourchoices[answer=C]
{质点 P 的运动方向向右}
{波的周期可能为 $0.27\Us$}
{波的频率可能为 $1.25\UHz$}
{波的传播速度可能为 $20\Ums$}

%% answer: C
%% memo:
\memoanswer{
本题原卷及材料中未提供详解，待补充。
}

\item
%% number: 4
%% paperName: 2005年广东理综第 4 题 4 分
%% typeId: 2
%% type: 多选题
%% chapter: 气体性质
%% point: 查理定律
%% method: 无
%% score: 4
%% degree: 650
%% duplicateId: 0
%% body:
封闭在气缸内一定质量的气体，如果保持气体体积不变，当温度升高时，以下说法正确的是\xzanswer{BD}

\fourchoices[answer=BD]
{气体的密度增大}
{气体的压强增大}
{气体分子的平均动能减小}
{每秒撞击单位面积器壁的气体分子数增多}

%% answer: BD
%% memo:
\memoanswer{
本题原卷及材料中未提供详解，待补充。
}

\item
%% number: 5
%% paperName: 2005年广东理综第 5 题 4 分
%% typeId: 1
%% type: 单选题
%% chapter: 光学
%% point: 光的折射
%% method: 无
%% score: 4
%% degree: 650
%% duplicateId: 0
%% body:
如图所示，一束白光通过玻璃棱镜发生色散现象，下列说法正确的是\xzanswer{B}

\onepicture[w=6.5cm]{figs/05.png}

\fourchoices[answer=B]
{红光的偏折最大，紫光的偏折最小}
{红光的偏折最小，紫光的偏折最大}
{玻璃对红光的折射率比紫光大}
{玻璃中紫光的传播速度比红光大}

%% answer: B
%% memo:
\memoanswer{
本题原卷及材料中未提供详解，待补充。
}

\item
%% number: 6
%% paperName: 2005年广东理综第 6 题 4 分
%% typeId: 2
%% type: 多选题
%% chapter: 电磁感应
%% point: 右手定则
%% method: 无
%% score: 4
%% degree: 450
%% duplicateId: 0
%% body:
如图所示，两根足够长的固定平行金属光滑导轨位于同一水平面，导轨上横放着两根相同的导体棒 ab、cd 与导轨构成矩形回路。导体棒的两端连接着处于压缩状态的两根轻质弹簧，两棒的中间用细线绑住，它们的电阻均为 $R$，回路上其余部分的电阻不计。在导轨平面内两导轨间有一竖直向下的匀强磁场。开始时，导体棒处于静止状态。剪断细线后，导体棒在运动过程中\xzanswer{AD}

\onepicture[w=5.5cm]{figs/06.png}

\fourchoices[answer=AD]
{回路中有感应电动势}
{两根导体棒所受安培力的方向相同}
{两根导体棒和弹簧构成的系统动量守恒，机械能守恒}
{两根导体棒和弹簧构成的系统动量守恒，机械能不守恒}

%% answer: AD
%% memo:
\memoanswer{
本题原卷及材料中未提供详解，待补充。
}

\item
%% number: 7
%% paperName: 2005年广东理综第 7 题 4 分
%% typeId: 2
%% type: 多选题
%% chapter: 光学
%% point: 光的折射
%% method: 无
%% score: 4
%% degree: 850
%% duplicateId: 0
%% body:
光纤通信是一种现代通信手段，它可以提供大容量、高速度、高质量的通信服务。目前，我国正在大力建设高质量的宽带光纤通信网络。下列说法正确的是\xzanswer{AD}

\fourchoices[answer=AD]
{光纤通信利用光作为载体来传播信息}
{光导纤维传递光信号是利用光的衍射原理}
{光导纤维传递光信号是利用光的色散原理}
{目前广泛应用的光导纤维是一种非常细的特制玻璃丝}

%% answer: AD
%% memo:
\memoanswer{
本题原卷及材料中未提供详解，待补充。
}

\item
%% number: 8
%% paperName: 2005年广东理综第 8 题 4 分
%% typeId: 2
%% type: 多选题
%% chapter: 光学
%% point: 光的电磁说
%% method: 无
%% score: 4
%% degree: 650
%% duplicateId: 0
%% body:
关于电磁场和电磁波，下列说法正确的是\xzanswer{ACD}

\fourchoices[answer=ACD]
{电磁波是横波}
{电磁波的传播需要介质}
{电磁波能产生干涉和衍射现象}
{电磁波中电场和磁场的方向处处相互垂直}

%% answer: ACD
%% memo:
\memoanswer{
本题原卷及材料中未提供详解，待补充。
}

\item
%% number: 9
%% paperName: 2005年广东理综第 9 题 4 分
%% typeId: 1
%% type: 单选题
%% chapter: 电磁感应
%% point: 法拉第电磁感应定律
%% method: 无
%% score: 4
%% degree: 650
%% duplicateId: 0
%% body:
钳形电流表的外形和结构如图（a）所示。图（a）中电流表的读数为 $1.2\UA$。图（b）中用同一电缆线绕了 3 匝，则\xzanswer{C}

\twopicture[wA=5.0cm, wB=2.4cm]{figs/09a.png}{figs/09b.png}

\fourchoices[answer=C]
{这种电流表能测直流电流，图（b）的读数为 $2.4\UA$}
{这种电流表能测交流电流，图（b）的读数为 $0.4\UA$}
{这种电流表能测交流电流，图（b）的读数为 $3.6\UA$}
{这种电流表既能测直流电流，又能测交流电流，图（b）的读数为 $3.6\UA$}

%% answer: C
%% memo:
\memoanswer{
本题原卷及材料中未提供详解，待补充。
}

\item
%% number: 10
%% paperName: 2005年广东理综第 10 题 4 分
%% typeId: 1
%% type: 单选题
%% chapter: 电场
%% point: 带电粒子的平衡
%% method: 无
%% score: 4
%% degree: 550
%% duplicateId: 0
%% body:
竖直放置的一对平行金属板的左极板上用绝缘线悬挂了一个带正电的小球，将平行金属板按图所示的电路图连接。绝缘线与左极板的夹角为 $\theta$。当滑动变阻器 $R$ 的滑片在 a 位置时，电流表的读数为 $I_{1}$，夹角为 $\theta_{1}$；当滑片在 b 位置时，电流表的读数为 $I_{2}$，夹角为 $\theta_{2}$，则\xzanswer{D}

\onepicture[w=6.5cm]{figs/10.png}

\fourchoices[answer=D]
{$\theta_{1}<\theta_{2}$，$I_{1}<I_{2}$}
{$\theta_{1}>\theta_{2}$，$I_{1}>I_{2}$}
{$\theta_{1}=\theta_{2}$，$I_{1}=I_{2}$}
{$\theta_{1}<\theta_{2}$，$I_{1}=I_{2}$}

%% answer: D
%% memo:
\memoanswer{
本题原卷及材料中未提供详解，待补充。
}

\item
%% number: 11
%% paperName: 2005年广东理综第 11 题 4 分
%% typeId: 4
%% type: 实验题
%% chapter: 直线运动
%% point: 打点计时器公式
%% method: 无
%% score: 4
%% degree: 550
%% duplicateId: 0
%% body:
\begin{enumerate}
	\item
	如图所示，在“用双缝干涉测光的波长”实验中，光具座上放置的光学元件依次为①光源、②\tkanswer[1.6cm]{滤光片}、③\tkanswer[1.6cm]{单缝}、④\tkanswer[1.6cm]{双缝}、⑤遮光筒、⑥光屏。对于某种单色光，为增加相邻亮纹（暗纹）间的距离，可采取\tkanswer[2.6cm]{减小双缝间距离}或\tkanswer[2.6cm]{增大双缝到屏的距离}的方法。

	\onepicture[w=10.5cm]{figs/11a.png}

	\item
	如图所示，某同学在做“研究匀变速直线运动”实验中，由打点计时器得到表示小车运动过程的一条清晰纸带，纸带上两相邻计数点的时间间隔为 $T=0.10\Us$，其中 $S_{1}=7.05\Ucm$、$S_{2}=7.68\Ucm$、$S_{3}=8.33\Ucm$、$S_{4}=8.95\Ucm$、$S_{5}=9.61\Ucm$、$S_{6}=10.26\Ucm$，则 A 点处瞬时速度的大小是 \tkanswer[1.4cm]{0.86} $\Ums$，小车运动的加速度计算表达式为 \tkanswer[4.5cm]{$a=\frac{(S_{4}+S_{5}+S_{6})-(S_{3}+S_{2}+S_{1})}{9T^{2}}$}，加速度的大小是 \tkanswer[1.4cm]{0.64} $\Umsq$（计算结果保留两位有效数字）。

	\onepicture[w=10.5cm]{figs/11b.png}
\end{enumerate}

%% answer: 
\jdanswer{
\begin{enumerate}
	\item
	②滤光片，③单缝，④双缝；减小双缝间距离，增大双缝到屏的距离。

	\item
	$0.86\Ums$，$a=\frac{(S_{4}+S_{5}+S_{6})-(S_{3}+S_{2}+S_{1})}{9T^{2}}$，$0.64\Umsq$
\end{enumerate}
}

%% memo:
\memoanswer{
本题原卷及材料中未提供详解，待补充。
}

\item
%% number: 12
%% paperName: 2005年广东理综第 12 题 11 分
%% typeId: 4
%% type: 实验题
%% chapter: 电路
%% point: 测电动势与内阻
%% method: 无
%% score: 11
%% degree: 450
%% duplicateId: 0
%% body:
热敏电阻是传感电路中常用的电子元件。现用伏安法研究热敏电阻在不同温度下的伏安特性曲线，要求特性曲线尽可能完整。已知常温下待测热敏电阻的阻值约 $4\sim5\UO$。热敏电阻和温度计插入带塞的保温杯中，杯内有一定量的冷水，其它备用的仪表和器具有：盛有热水的热水杯（图中未画出）、电源（$3\UV$、内阻可忽略）、直流电流表（内阻约 $1\UO$）、直流电压表（内阻约 $5\UkO$）、滑动变阻器（$0\sim20\UO$）、开关、导线若干。

\begin{enumerate}
	\item
	在图（a）的方框中画出实验电路图，要求测量误差尽可能小。

	\drawpicanswer{\onepicture[w=5.0cm]{figs/12a.png}}{\onepicture[w=5.0cm]{figs/12a_an.png}}

	\item
	根据电路图，在图（b）的实物图上连线。

	\drawpicanswer{\onepicture[w=8.0cm]{figs/12b.png}}{\onepicture[w=8.0cm]{figs/12b_an.png}}

	\item
	简要写出完成接线后的主要实验步骤\tkanswer[4.0cm]{往保温杯中加入一些热水，待温度稳定时读出温度计值；调节滑动变阻器，快速测出几组电流表和电压表的值；重复操作，测量不同温度下的数据；绘出各测量温度下热敏电阻的伏安特性曲线}。
\end{enumerate}

%% answer: 
\jdanswer{
\begin{enumerate}
	\item
	外接，分压，电流表用 $3\UA$ 量程，电压表用 $3\UV$ 量程。

	\item
	实物连线见所述电路图。

	\item
	①往保温杯中加入一些热水，待温度稳定时读出温度计值；②调节滑动变阻器，快速测出几组电流表和电压表的值；③重复①～②，测量不同温度下的数据；④绘出各测量温度下热敏电阻的伏安特性曲线。
\end{enumerate}
}

%% memo:
\memoanswer{
本题原卷及材料中未提供详解，待补充。
}

\item
%% number: 13
%% paperName: 2005年广东理综第 13 题 8 分
%% typeId: 3
%% type: 填空题
%% chapter: 气体性质
%% point: 热力学第一定律
%% method: 无
%% score: 8
%% degree: 650
%% duplicateId: 0
%% body:
\begin{enumerate}
	\item
	如图所示，氢原子从 $n>2$ 的某一能级跃迁到 $n=2$ 的能级，辐射出能量为 $2.55\UeV$ 的光子。问最少要给基态的氢原子提供多少电子伏特的能量，才能使它辐射上述能量的光子？请在图中画出获得该能量后的氢原子可能的辐射跃迁图。

	\onepicture[align=right]{figs/13a.png}

	\item
	热力学第二定律常见的表述有两种。第一种表述：不可能使热量由低温物体传递到高温物体，而不引起其他变化；第二种表述：不可能从单一热源吸收热量并把它全部用来做功，而不引起其他变化。图（a）是根据热力学第二定律的第一种表述画出的示意图：外界对制冷机做功，使热量从低温物体传递到高温物体。请你根据第二种表述完成示意图（b）。根据你的理解，热力学第二定律的实质是\tkanswer[4.0cm]{自然界中进行的与热现象有关的宏观物理过程都具有方向性}。

	\onepicture[num=a, align=right, w=4.0cm]{figs/13b.png}

	\drawpicanswer{\onepicture[num=b, align=right, w=4.0cm]{figs/13c.png}}{\onepicture[num=b, align=right, w=4.0cm]{figs/13c_an.png}}
\end{enumerate}

%% answer: 
\jdanswer{
\begin{enumerate}
	\item
	$12.75\UeV$，略（共 6 条）。

	\item
	自然界中进行的与热现象有关的宏观物理过程都具有方向性。
\end{enumerate}
}

%% memo:
\memoanswer{
本题原卷及材料中未提供详解，待补充。
}

\item
%% number: 14
%% paperName: 2005年广东理综第 14 题 12 分
%% typeId: 6
%% type: 计算题
%% chapter: 功和能
%% point: 机械能守恒定律
%% method: 无
%% score: 12
%% degree: 650
%% duplicateId: 0
%% body:
如图所示，半径 $R=0.40\Um$ 的光滑半圆环轨道处于竖直平面内，半圆环与粗糙的水平地面相切于圆环的端点 A。一质量 $m=0.10\Ukg$ 的小球，以初速度 $v_{0}=7.0\Ums$ 在水平地面上向左作加速度 $a=3.0\Umsq$ 的匀减速直线运动，运动 $4.0\Um$ 后，冲上竖直半圆环，最后小球落在 C 点。求 A、C 间的距离（取重力加速度 $g=10\Umsq$）。

\onepicture[align=right]{figs/14.png}

%% answer: 
\jdanswer{$1.2\Um$}

%% memo:
\memoanswer{
本题原卷及材料中未提供详解，待补充。
}

\item
%% number: 15
%% paperName: 2005年广东理综第 15 题 13 分
%% typeId: 6
%% type: 计算题
%% chapter: 曲线运动
%% point: 人造地球卫星
%% method: 无
%% score: 13
%% degree: 450
%% duplicateId: 0
%% body:
已知万有引力常量 $G$，地球半径 $R$，月球和地球之间的距离 $r$，同步卫星距地面的高度 $h$，月球绕地球的运转周期 $T_{1}$，地球的自转周期 $T_{2}$，地球表面的重力加速度 $g$。某同学根据以上条件，提出一种估算地球质量 $M$ 的方法：

同步卫星绕地球作圆周运动，由 $\frac{GMm}{h^{2}}=m\left(\frac{2\pi}{T}\right)^{2}h$，得 $M=\frac{4\pi^{2}h^{3}}{GT^{2}}$

\begin{enumerate}
	\item
	请判断上面的结果是否正确，并说明理由。如不正确，请给出正确的解法和结果。
	\item
	请根据已知条件再提出两种估算地球质量的方法并解得结果。
\end{enumerate}

%% answer: 
\jdanswer{
\begin{enumerate}
	\item
	上面的结果是错误的，地球半径 $R$ 在计算过程中不能忽略。

	$M=\frac{4\pi^{2}(R+h)^{3}}{GT_{2}^{2}}$

	\item
	$M=\frac{gR^{2}}{G}$；$M=\frac{4\pi^{2}r^{3}}{GT_{1}^{2}}$
\end{enumerate}
}

%% memo:
\memoanswer{
（1）不正确。应为 $G\frac{Mm}{(R+h)^{2}}=m\left(\frac{2\pi}{T_{2}}\right)^{2}(R+h)$，得 $M=\frac{4\pi^{2}(R+h)^{3}}{GT_{2}^{2}}$。

（2）由 $mg=\frac{GMm}{R^{2}}$，得 $M=\frac{gR^{2}}{G}$；或由 $\frac{GMm}{r^{2}}=m\left(\frac{2\pi}{T_{1}}\right)^{2}r$，得 $M=\frac{4\pi^{2}r^{3}}{GT_{1}^{2}}$。
}

\item
%% number: 16
%% paperName: 2005年广东理综第 16 题 16 分
%% typeId: 6
%% type: 计算题
%% chapter: 磁场
%% point: 带电粒子在磁场中的运动
%% method: 无
%% score: 16
%% degree: 450
%% duplicateId: 0
%% body:
如图所示，在一个圆形区域内，两个方向相反且都垂直于纸面的匀强磁场分布在以直径 $A_{2}A_{4}$ 为边界的两个半圆形区域Ⅰ、Ⅱ中，$A_{2}A_{4}$ 与 $A_{1}A_{3}$ 的夹角为 $60^{\circ}$。一质量为 $m$、带电量为 $+q$ 的粒子以某一速度从Ⅰ区的边缘点 $A_{1}$ 处沿与 $A_{1}A_{3}$ 成 $30^{\circ}$ 角的方向射入磁场，随后该粒子以垂直于 $A_{2}A_{4}$ 的方向经过圆心 $O$ 进入Ⅱ区，最后再从 $A_{4}$ 处射出磁场。已知该粒子从射入到射出磁场所用的时间为 $t$，求Ⅰ区和Ⅱ区中磁感应强度的大小（忽略粒子重力）。

\onepicture[align=right]{figs/16.png}

%% answer: 
\jdanswer{$B_{1}=\frac{4\pi m}{3qt}$，$B_{2}=\frac{8\pi m}{3qt}$}

%% memo:
\memoanswer{
本题原卷及材料中未提供详解，待补充。
}

\item
%% number: 17
%% paperName: 2005年广东理综第 17 题 16 分
%% typeId: 6
%% type: 计算题
%% chapter: 磁场
%% point: 带电粒子在复合场中的运动
%% method: 无
%% score: 16
%% degree: 450
%% duplicateId: 0
%% body:
如图所示，一半径为 $r$ 的圆形导线框内有一匀强磁场，磁场方向垂直于导线框所在平面，导线框的左端通过导线接一对水平放置的平行金属板，两板间的距离为 $d$，板长为 $l$，$t=0$ 时，磁场的磁感应强度 $B$ 从 $B_{0}$ 开始均匀增大，同时，在板 2 的左端且非常靠近板 2 的位置有一质量为 $m$、带电量为 $-q$ 的液滴以初速度 $v_{0}$ 水平向右射入两板间，该液滴可视为质点。

\begin{enumerate}
	\item
	要使该液滴能从两板间射出，磁感应强度随时间的变化率 $k$ 应满足什么条件？
	\item
	要使该液滴能从两板间右端的中点射出，磁感应强度 $B$ 与时间 $t$ 应满足什么关系？
\end{enumerate}

\onepicture[align=right]{figs/17.png}

%% answer: 
\jdanswer{
\begin{enumerate}
	\item
	$\frac{mgd}{q\pi r^{2}}\leqslant k\leqslant\frac{mgd}{q\pi r^{2}}+\frac{2mv_{0}^{2}d^{2}}{\pi r^{2}ql^{2}}$

	\item
	$B=B_{0}+\left(\frac{mgd}{q\pi r^{2}}+\frac{mv_{0}^{2}d^{2}}{\pi r^{2}ql^{2}}\right)$
\end{enumerate}
}

%% memo:
\memoanswer{
本题原卷及材料中未提供详解，待补充。
}

\item
%% number: 18
%% paperName: 2005年广东理综第 18 题 17 分
%% typeId: 6
%% type: 计算题
%% chapter: 运动和力
%% point: 动量和能量
%% method: 无
%% score: 17
%% degree: 350
%% duplicateId: 0
%% body:
如图所示，两个完全相同的质量为 $m$ 的木板 A、B 置于水平地面上，它们的间距 $s=2.88\Um$。质量为 $2m$，大小可忽略的物块 C 置于 A 板的左端。C 与 A 之间的动摩擦因数为 $\mu_{1}=0.22$，A、B 与水平地面之间的动摩擦因数为 $\mu_{2}=0.10$，最大静摩擦力可以认为等于滑动摩擦力。开始时，三个物体处于静止状态。现给 C 施加一个水平向右，大小为 $\frac{2}{5}mg$ 的恒力 $F$，假定木板 A、B 碰撞时间极短且碰撞后粘连在一起，要使 C 最终不脱离木板，每块木板的长度至少应为多少？

\onepicture[align=right]{figs/18.png}

%% answer: 
\jdanswer{$l=3\Um$}

%% memo:
\memoanswer{
第一阶段拉力 $F$ 小于 CA 间最大静摩擦力，因此 CA 共同加速到与 B 相碰，该过程对 CA 用动能定理：

$F-\mu_{2}\cdot 3mgs=\frac{3mv_{1}^{2}}{2}$，得 $v_{1}=8\sqrt{0.3}\Ums$。

AB 相碰瞬间，AB 动量守恒，碰后共同速度 $v_{2}=4\sqrt{0.3}\Ums$。

C 在 AB 上滑行全过程，ABC 系统所受合外力为零，动量守恒，C 到 B 右端时恰好达到共速：

$2mv_{1}+2mv_{2}=4mv$，因此共同速度 $v=6\sqrt{0.3}\Ums$。

C 在 AB 上滑行全过程用能量守恒：$F\cdot 2L=\frac{4mv^{2}}{2}-\left(\frac{2mv_{1}^{2}}{2}+\frac{2mv_{2}^{2}}{2}\right)+\mu_{1}\cdot 2mg\cdot 2L$，

得 $l=3\Um$。
}

\end{enumerate}

\end{document}
